% !TeX root = paper.tex
\newcommand*{\PathToOutput}{../_output}

\documentclass[12pt]{article}

\usepackage{my_article_header}
\usepackage{my_common_header}

\title{Replication of Jiang et al. (2024):\\
Monetary Tightening and U.S. Bank Fragility}

\author{
Anoushka Gehani \\
Hashir Bawany
}

\date{\today}

\begin{document}

\maketitle

\doublespacing

\begin{abstract}

This project replicates key results from Jiang, Matvos, Piskorski, and Seru (2024), 
which studies the impact of rising interest rates on the U.S. banking system. 
Using regulatory bank call report data and market price indices, the project 
reconstructs measures of bank balance sheet exposure to interest rate risk. 
The analysis reproduces several tables and figures from the original paper 
and extends them using updated data through 2025. The results confirm that 
rapid monetary tightening significantly reduced the market value of bank assets 
and that banks with high shares of uninsured deposits were more vulnerable to 
deposit runs.

\end{abstract}

\section{Introduction}

In 2022 and 2023 the Federal Reserve implemented one of the fastest monetary 
tightening cycles in recent decades in response to rising inflation. 
Rapid increases in interest rates reduced the market value of long-duration 
assets such as Treasury securities and mortgage-backed securities held by banks.

Jiang et al. (2024) study how these valuation losses affected the stability 
of the U.S. banking system. Their analysis highlights the interaction between 
unrealized losses on bank balance sheets and the funding structure of banks, 
particularly the share of uninsured deposits.

This project replicates several key results from the paper and updates the 
analysis using more recent regulatory data.

\section{Data}

The primary dataset used in this project is the U.S. bank Call Report, a 
quarterly regulatory filing submitted by all insured commercial banks to 
federal regulators through the Federal Financial Institutions Examination 
Council (FFIEC).

The Call Reports provide detailed balance sheet information including:

\begin{itemize}
\item total assets and liabilities
\item loan portfolios
\item securities holdings
\item deposit composition
\end{itemize}

To estimate mark-to-market losses, the project combines these regulatory data 
with market price series including:

\begin{itemize}
\item U.S. Treasury yield curves
\item mortgage-backed securities price indices
\item macroeconomic data from the FRED database
\end{itemize}

All data are collected and processed using an automated Python pipeline.

\section{Replication of Paper Results}

This section reproduces the main statistics reported in Jiang et al. (2024). 
The code reconstructs the bank-level exposure to interest rate shocks and 
aggregates these exposures across banks.

\begin{table}[h]
\centering
\caption{Replication of Table 1}
\input{\PathToOutput/table1.tex}
\caption*{Estimated mark-to-market losses across different categories of banks.}
\end{table}

\section{Updated Results Through 2025}

Using updated Call Report data and market price series, the analysis extends 
the calculations in the original paper through 2025.

\begin{table}[h]
\centering
\caption{Updated Market Value Losses}
\input{\PathToOutput/table1_updated.tex}
\caption*{Updated estimates of unrealized losses using data through 2025.}
\end{table}

\section{Summary Statistics}

To better understand the underlying data, this section provides descriptive 
statistics on bank balance sheets and asset exposures.

\begin{table}[h]
\centering
\caption{Summary Statistics of Bank Balance Sheets}
\input{\PathToOutput/summary_stats.tex}
\caption*{Summary statistics of bank assets, securities holdings, and deposit composition.}
\end{table}

\begin{figure}[h]
\centering
\includegraphics[width=0.8\linewidth]{\PathToOutput/bank_asset_distribution.png}
\caption{Distribution of Bank Asset Exposure}
\end{figure}

\section{Discussion}

The replication successfully reproduces several key results from the original 
paper. The results confirm that increases in interest rates led to substantial 
declines in the market value of bank assets.

Updating the analysis through 2025 shows that exposure to interest rate risk 
remains significant across many institutions. Banks with higher shares of 
uninsured deposits appear more vulnerable to rapid deposit withdrawals.

The main challenges in this project involved collecting and cleaning the 
regulatory call report data and ensuring that the calculations matched the 
methodology used in the original paper.

\section{Conclusion}

This project demonstrates how reproducible analytical pipelines can be used to 
replicate complex empirical results in finance. By automating the data collection, 
cleaning, and analysis steps, the project produces updated estimates of bank 
fragility that can be easily extended as new data become available.

\bibliographystyle{jpe}
\bibliography{bibliography.bib}

\end{document}
